% ==============================================================
% Differenz-Relation-Struktur (DRS)
% Form als fundamentaler ontologischer Begriff
% Version 0.3 (Entwurf)
% Teil 1
% Kompiliert mit: pdfLaTeX
% ==============================================================

\documentclass[11pt,a4paper]{article}

% ---------- Grundpakete ----------
\usepackage[utf8]{inputenc}
\usepackage[T1]{fontenc}
\usepackage[ngerman,provide=*]{babel}

\usepackage{amsmath,amssymb,amsthm}
\usepackage{mathtools}

\usepackage{geometry}
\usepackage{parskip}
\usepackage{enumitem}
\usepackage{booktabs}
\usepackage{xcolor}
\usepackage{hyperref}

% ---------- Layout ----------
\geometry{margin=2.5cm}

\setlist{
itemsep=0.25em,
topsep=0.25em
}

\hypersetup{
colorlinks=true,
linkcolor=blue,
urlcolor=blue,
pdftitle={Differenz-Relation-Struktur (DRS)},
pdfauthor={Forschungsprogramm Ontophänomenalismus}
}

% ---------- Theorem ----------
\theoremstyle{plain}
\newtheorem{theorem}{Theorem}[section]
\newtheorem{proposition}{Proposition}[section]
\newtheorem{lemma}{Lemma}[section]

\theoremstyle{definition}
\newtheorem{definition}{Definition}[section]
\newtheorem{remark}{Bemerkung}[section]
\newtheorem{methodennote}{Methodennote}[section]
\newtheorem{hypothese}{Arbeitshypothese}[section]

% ==============================================================
% Titel
% ==============================================================

\title{
\textbf{Differenz--Relation--Struktur (DRS)}\\[0.4em]
{\large Form als fundamentaler ontologischer Begriff}\\[0.4em]
{\normalsize Version 0.3 (Entwurf)}
}

\author{
Forschungsprogramm Ontophänomenalismus\\
\small Siehe TRANSPARENCY für methodologische Hinweise.
}

\date{Juli 2026}

\begin{document}

\maketitle

\begin{abstract}

Dieses Dokument untersucht die Hypothese, dass der Begriff der
\emph{Form} den allgemeinsten ontologischen Grundbegriff des
Ontophänomenalismus darstellt.

Während die bisherigen Dokumente des OP Begriffe wie
Existenz,
Bestimmtheit,
Relation,
Information
und Struktur teilweise getrennt behandeln,
wird hier untersucht,
ob sich diese Begriffe auf ein gemeinsames Fundament zurückführen lassen.

Die leitende Vermutung lautet:

\begin{center}

\emph{Existenz ist nur als Form möglich.}

\end{center}

Form wird dabei nicht geometrisch verstanden,
sondern als stabile Organisation von Differenzen und Relationen.
Das Dokument besitzt den Status einer Arbeitshypothese
(MKO Ebene 4--5)
und dient der Vorbereitung einer späteren mathematischen
Formalisierung innerhalb der MGO und PST.

\end{abstract}

\tableofcontents

\newpage

% ==============================================================
\section{Einordnung}
% ==============================================================

\begin{methodennote}[Stellung innerhalb des OP]

DRS besitzt gegenwärtig den Charakter eines Forschungsdokuments.

Es erweitert nicht die Axiome des Ontophänomenalismus,
sondern untersucht,
ob mehrere bereits verwendete Grundbegriffe
unter einem gemeinsamen Prinzip zusammengeführt werden können.

Die wichtigsten Schnittstellen sind:

\begin{itemize}

\item
OP:
Existenz,
Bestimmtheit,
Unterscheidbarkeit.

\item
FSO:
Formale Semantik,
Diskursdomänen,
Referenz.

\item
MGO:
Mathematische Beschreibung von Struktur,
Operatoren,
Kategorien
und Projektionen.

\item
PST:
Entstehung,
Transformation
und Dynamik von Formen.

\item
PGT:
Selbstreferenzielle Formen,
Bewusstsein,
Perspektiven.

\end{itemize}

\end{methodennote}

% ==============================================================
\section{Motivation}
% ==============================================================

Der Ontophänomenalismus verwendet bereits mehrere eng miteinander
verbundene Grundbegriffe.

Unter anderem:

\begin{itemize}

\item Existenz

\item Bestimmtheit

\item Differenz

\item Relation

\item Struktur

\item Information

\end{itemize}

Bislang werden diese Begriffe jedoch weitgehend unabhängig
voneinander eingeführt.

Dies wirft die Frage auf,
ob sie möglicherweise verschiedene Erscheinungsformen
eines allgemeineren Prinzips darstellen.

Die leitende Fragestellung lautet deshalb:

\begin{quote}

Ist Form die allgemeinste Beschreibung dessen,
was minimale Bestimmtheit überhaupt bedeutet?

\end{quote}

Falls diese Vermutung zutrifft,
ergäbe sich eine erhebliche Vereinfachung
der ontologischen Grundstruktur.

Existenz müsste dann nicht mehr unmittelbar
über Bestimmtheit definiert werden,
sondern Bestimmtheit selbst wäre lediglich
eine Eigenschaft von Form.

Ebenso würden Information,
Struktur,
Relation
und möglicherweise sogar Raum
nicht mehr als unabhängige Begriffe erscheinen,
sondern als verschiedene Aspekte derselben
formalen Grundlage.

% ==============================================================
\section{Ausgangspunkt}
% ==============================================================

Die Überlegungen bauen unmittelbar auf drei bereits eingeführten
Prinzipien auf.

\subsection*{Axiom 1}

Etwas existiert.

\subsection*{Unterscheidbarkeits-Prämisse (UP)}

Existenz setzt minimale Bestimmtheit voraus.

Was prinzipiell von nichts unterschieden werden kann,
besitzt mindestens eine Eigenschaft.

\subsection*{Konsequenz}

Bestimmtheit verlangt immer mindestens
eine unterscheidbare Struktur.

Die hier entwickelte Arbeitshypothese lautet daher:

\[
\boxed{
\text{Bestimmtheit}
\Longleftrightarrow
\text{Form}
}
\]

Diese Identifikation ist gegenwärtig noch keine Aussage des OP,
sondern Gegenstand der vorliegenden Untersuchung.

% ==============================================================
\section{Differenz als primitiver Begriff}
% ==============================================================

\begin{definition}[Differenz]

Eine Differenz bezeichnet die elementare Möglichkeit,
zwei Zustände,
Eigenschaften
oder Relationen voneinander zu unterscheiden.

Formal schreiben wir

\[
a \neq b.
\]

Dabei beschreibt die Differenz noch keine Richtung,
keine Entfernung
und keine Metrik.

Sie bezeichnet ausschließlich Nichtidentität.

\end{definition}

\begin{remark}

Differenz wird hier nicht weiter definiert.

Sie besitzt den Status eines primitiven Begriffs,
vergleichbar mit Punkt oder Menge
in axiomatischen Theorien.

\end{remark}

% ==============================================================
\section{Relation}
% ==============================================================

\begin{definition}[Relation]

Eine Relation beschreibt eine strukturierte Verknüpfung
zwischen mindestens zwei unterscheidbaren Elementen.

Formal

\[
R
\subseteq
D\times D.
\]

\end{definition}

\begin{remark}

Eine Relation setzt Differenzen voraus.

Sind keinerlei unterscheidbare Pole vorhanden,
reduziert sich jede Relation
auf reine Identität.

\end{remark}

\begin{proposition}[Kooriginarität]

Differenz und Relation sind kooriginär.

Differenzen erzeugen Relationen.

Relationen erhalten Differenzen.

Keine der beiden Strukturen besitzt
innerhalb dieser Arbeitshypothese
einen absoluten Vorrang.

\end{proposition}

Dies lässt sich symbolisch schreiben als

\[
\boxed{
\text{Differenz}
\Longleftrightarrow
\text{Relation}
}
\]

Diese Aussage ist ausdrücklich keine logische Äquivalenz,
sondern beschreibt gegenseitige ontologische Abhängigkeit.

% ==============================================================
\section{Form}
% ==============================================================

\begin{definition}[Form]

Eine Form ist eine stabile Organisation
von Differenzen und Relationen.

Formal schreiben wir

\[
F=(D,R,I).
\]

Dabei bezeichnet

\begin{itemize}

\item
\(D\)
die Menge der Differenzen,

\item
\(R\)
die Relationen zwischen diesen Differenzen,

\item
\(I\)
eine Invarianzbedingung,
welche die Identität der Form
unter zulässigen Transformationen beschreibt.

\end{itemize}

\end{definition}

\begin{remark}

Die Invarianz ersetzt bewusst den bisherigen Begriff
der Stabilität.

Stabilität beschreibt lediglich Verhalten.

Invarianz beschreibt dagegen,
welche Eigenschaften trotz Transformation
erhalten bleiben.

Damit wird Form unabhängig
von ihrem konkreten Träger.

\end{remark}

% ==============================================================
\section{Erstes Formprinzip}
% ==============================================================

\begin{proposition}[Formprinzip]

Existenz setzt Form voraus.

Formal

\[
\boxed{
\text{Existenz}
\Longrightarrow
\text{Form}
}
\]

\end{proposition}

\begin{proof}[Konzeptuelle Herleitung]

Existenz verlangt nach UP
mindestens eine Bestimmtheit.

Bestimmtheit verlangt mindestens
eine Differenz.

Differenzen erzeugen Relationen.

Stabile Relationen bilden eine Form.

Damit besitzt jede Existenz
mindestens eine Form.

\end{proof}

\begin{remark}

Die Umkehrung

\[
\text{Form}
\Longrightarrow
\text{Existenz}
\]

wird bewusst noch nicht behauptet.

Ob jede mathematisch beschreibbare Form
ontologische Existenz besitzt,
ist eine offene Forschungsfrage
und berührt unmittelbar
den Status mathematischer Objekte.

\end{remark}

% ==============================================================
% Ende Teil 1
% ==============================================================

% Fortsetzung in Teil 2

% ==============================================================
% DRS – Differenz-Relation-Struktur
% Version 0.3 (Entwurf)
% Teil 2
% ==============================================================

% ==============================================================
\section{Form und Information}
% ==============================================================

\begin{methodennote}[Leitidee]

In der Literatur wird Information häufig als fundamentaler Begriff
behandelt.

Die hier entwickelte Arbeitshypothese verfolgt die umgekehrte Richtung:

\begin{center}
Information setzt Form voraus.
\end{center}

Form besitzt daher im DRS den ontologischen Vorrang.

\end{methodennote}

%--------------------------------------------------------------

\begin{definition}[Information]

Information ist eine Interpretation oder Auslesung einer Form
relativ zu einem Beobachter,
einem Referenzsystem
oder einer Projektion.

Formal schreiben wir

\[
\mathcal I = \Phi(F)
\]

wobei

\[
\Phi
\]

einen Interpretationsoperator bezeichnet.

\end{definition}

\begin{remark}

Dies trennt erstmals drei Begriffe sauber voneinander.

\[
\boxed{
\text{Form}
\neq
\text{Information}
\neq
\text{Bedeutung}
}
\]

Form existiert unabhängig davon,
ob sie interpretiert wird.

Information entsteht erst durch Interpretation.

Bedeutung setzt zusätzlich einen semantischen Kontext voraus
(FSO).

\end{remark}

% ==============================================================
\section{Form und Darstellung}
% ==============================================================

Verschiedene Darstellungen können dieselbe Form besitzen.

Die Form ist daher trägerunabhängig.

Formal:

\[
F
=
[\rho_1]
=
[\rho_2]
=
\cdots
=
[\rho_n]
\]

wobei

\[
[\rho]
\]

eine Äquivalenzklasse zulässiger Darstellungen bezeichnet.

Damit besitzt Form eine deutlich höhere Allgemeinheit
als jede konkrete Repräsentation.

%--------------------------------------------------------------

\begin{remark}

Eine Melodie kann

\begin{itemize}

\item gespielt,

\item gesungen,

\item als Notenschrift,

\item digital codiert

\end{itemize}

werden.

Alle diese Repräsentationen unterscheiden sich materiell,
realisieren jedoch dieselbe Form.

\end{remark}

% ==============================================================
\section{Das Kompressionsprinzip}
% ==============================================================

Die Unterscheidung zwischen Form und Darstellung
führt unmittelbar auf ein allgemeines Kompressionsprinzip.

\begin{proposition}[Kompressionsprinzip]

Eine Form kann wesentlich einfacher sein
als ihre vollständige Entfaltung.

\end{proposition}

\begin{remark}

Die Kreiszahl

\[
\pi
\]

liefert ein anschauliches Beispiel.

Ihre Dezimalentwicklung besitzt unendlich viele Stellen.

Die zugrunde liegende Form

\[
\pi
=
\frac{\text{Umfang}}{\text{Durchmesser}}
\]

ist dagegen äußerst kompakt.

Unendliche Darstellung impliziert daher
keine komplexe Form.

\end{remark}

Dieses Prinzip dürfte später
für Attraktoren,
Projektionsoperatoren
und die Beschreibung physikalischer Zustände
von erheblicher Bedeutung sein.

% ==============================================================
\section{Transformationen von Formen}
% ==============================================================

Formen sind nicht statisch.

Sie können ineinander übergehen,
ohne ihre Identität vollständig zu verlieren.

\begin{definition}[Transformation]

Eine Transformation ist ein Operator

\[
T :
F
\rightarrow
F'
\]

der eine Form
in eine andere Form überführt.

\end{definition}

Transformationen können beispielsweise sein

\begin{itemize}

\item Projektionen,

\item Symmetrien,

\item Deformationen,

\item Kompositionen,

\item Zerlegungen.

\end{itemize}

%--------------------------------------------------------------

\begin{definition}[Invariant]

Eine Eigenschaft

\[
P(F)
\]

heißt invariant bezüglich

\[
T,
\]

wenn gilt

\[
P(F)
=
P(T(F)).
\]

\end{definition}

Invarianten definieren die Identität einer Form.

Nicht alle Eigenschaften müssen erhalten bleiben.

Entscheidend ist,
welche Eigenschaften unter zulässigen Transformationen
konstant bleiben.

% ==============================================================
\section{Minimalformprinzip}
% ==============================================================

Die bisher entwickelten Begriffe erlauben
eine stärkere Formulierung.

\begin{proposition}[Minimalformprinzip]

Jede Existenz besitzt mindestens eine minimale Form.

Formal

\[
\forall x
\,
\big(
E(x)
\rightarrow
\exists F_x
\big).
\]

\end{proposition}

\begin{remark}

Das Minimalformprinzip erweitert unmittelbar
die Unterscheidbarkeits-Prämisse.

Existenz verlangt Bestimmtheit.

Bestimmtheit verlangt Form.

Form verlangt mindestens eine Differenz
und mindestens eine Relation.

\end{remark}

% ==============================================================
\section{Raum als Form}
% ==============================================================

Im Ontophänomenalismus
wird Raum nicht als primitive Entität verstanden.

Vielmehr entsteht Raum
durch stabile relationale Formen.

\begin{hypothese}[Emergenz des Raumes]

Raum ist eine emergente Form
eines Netzes relationaler Möglichkeiten.

\end{hypothese}

Damit besitzt Raum
keinen ontologischen Vorrang.

Vorrang besitzt ausschließlich
die zugrunde liegende Form.

Dies stimmt mit der projektiven Interpretation
des OPP überein.

% ==============================================================
\section{Entropie}
% ==============================================================

Auch Entropie erscheint
unter diesem Blickwinkel
nicht mehr fundamental.

\begin{remark}

Sowohl

\begin{itemize}

\item Shannon-Entropie,

\item Von-Neumann-Entropie,

\item thermodynamische Entropie

\end{itemize}

setzen einen strukturierten Zustandsraum voraus.

Ein Zustandsraum
ist jedoch bereits eine Form.

Ohne Form
existieren weder Zustände
noch Wahrscheinlichkeiten.

Entropie setzt daher Form voraus.

\end{remark}

Formal:

\[
\boxed{
\text{Form}
\Longrightarrow
\text{Zustandsraum}
\Longrightarrow
\text{Entropie}
}
\]

Die Umkehrung gilt im Allgemeinen nicht.

% ==============================================================
\section{Möglichkeiten}
% ==============================================================

Auch Möglichkeiten
lassen sich formtheoretisch interpretieren.

\begin{definition}[Möglichkeitsraum]

Ein Möglichkeitsraum
ist eine Menge potentiell realisierbarer Formen.

\[
\Omega
=
\{F_i\}
\]

\end{definition}

Physikalische Realität
wählt innerhalb des OP
nicht aus formlosen Möglichkeiten,
sondern aus strukturierten Formen.

Dies verbindet DRS
direkt mit der Projektiven Strukturtheorie.

% ==============================================================
\section{Zwischenfazit}
% ==============================================================

Die bisherige Argumentation lässt sich
in einer ersten Hierarchie zusammenfassen.

\[
\boxed{
\begin{aligned}
&
\text{Differenz}
\\[0.3em]
\Longleftrightarrow\;&
\text{Relation}
\\[0.3em]
\Longrightarrow\;&
\text{Form}
\\[0.3em]
\Longrightarrow\;&
\text{Bestimmtheit}
\\[0.3em]
\Longrightarrow\;&
\text{Existenz}
\end{aligned}
}
\]

Die genaue Richtung einzelner Implikationen
bleibt Gegenstand weiterer Untersuchungen.

Insbesondere ist noch offen,

\[
\text{Form}
\Longrightarrow
\text{Existenz}
\]

oder

\[
\text{Existenz}
\Longleftrightarrow
\text{Form}
\]

ob eine dieser stärkeren Aussagen
innerhalb des OP begründet werden kann.

% ==============================================================
% Ende Teil 2
% ==============================================================

% Fortsetzung in Teil 3
% ==============================================================
% DRS – Differenz-Relation-Struktur
% Version 0.3 (Entwurf)
% Teil 3
% ==============================================================

%==============================================================
\section{Bewusstsein als Form}
%==============================================================

\begin{methodennote}[Status]

Die folgenden Überlegungen stellen eine offene Arbeitshypothese (MKO Ebene 5)
dar. Sie dienen als mögliche Verbindung zwischen DRS und der Projektiven
Gefälletheorie (PGT).

\end{methodennote}

Im bisherigen OP wird Bewusstsein nicht als eigenständige Substanz verstanden,
sondern als emergente Eigenschaft bestimmter relationaler Strukturen.

Die DRS-Hypothese erlaubt eine allgemeinere Formulierung.

\begin{hypothese}[Bewusstsein]

Bewusstsein ist eine besondere Klasse
selbstreferenzieller und dynamisch stabiler Formen.

\end{hypothese}

Selbstreferenz bedeutet dabei nicht lediglich Rekursion, sondern die Fähigkeit
einer Form, ihre eigene Struktur in ihre weitere Entwicklung einzubeziehen.

Formal könnte dies durch einen Operator

\[
\Sigma(F)
\]

beschrieben werden, für den gilt

\[
\Sigma(F)\subseteq F
\]

oder allgemeiner

\[
F \rightarrow \Sigma(F) \rightarrow F'.
\]

Ob diese Selbstreferenz bereits hinreichend für Bewusstsein ist oder lediglich
eine notwendige Bedingung darstellt, bleibt offen.

%==============================================================
\section{Projektive Interpretation}
%==============================================================

Die DRS lässt sich als strukturelle Grundlage der Projektiven Strukturtheorie
(PST) interpretieren.

\begin{hypothese}

Projektionsoperatoren erzeugen keine Form aus Formlosigkeit.

Sie bilden bestehende Formen in neue Formen ab.

\[
P :
F_{OPP}
\longrightarrow
F_{phys}
\]

\end{hypothese}

Damit verschiebt sich der Fokus des OP:

Nicht Materie ist fundamental,
sondern Form.

Nicht Raum ist fundamental,
sondern Form.

Nicht Information ist fundamental,
sondern Form.

Die physikalische Welt erscheint dann als eine Familie projektierter Formen.

%==============================================================
\section{Form und mathematische Struktur}
%==============================================================

Eine zentrale offene Frage betrifft die geeignete mathematische Beschreibung
von Formen.

Derzeit erscheinen mehrere Kandidaten geeignet.

\begin{enumerate}

\item Graphentheorie

\item Hypergraphen

\item Kategorien und Funktoren

\item Topologie

\item Operatoralgebren

\item Dynamische Systeme

\end{enumerate}

Jeder dieser Ansätze beschreibt unterschiedliche Aspekte von Form.

Eine zukünftige MGO-Version sollte untersuchen,
welcher Formalismus die größte Ausdrucksstärke bei möglichst geringer
Komplexität besitzt.

Insbesondere erscheint die Kategorientheorie vielversprechend,
da sie Beziehungen zwischen Strukturen statt einzelner Objekte beschreibt.

%==============================================================
\section{Invarianz als Identitätsprinzip}
%==============================================================

Die bisherige Definition

\[
F=(D,R,S)
\]

liefert eine operative Beschreibung einer Form.

Möglicherweise lässt sich der Begriff jedoch noch allgemeiner fassen.

\begin{hypothese}[Form als Invarianzklasse]

Eine Form ist die Äquivalenzklasse aller Differenz-Relation-Strukturen,
die bezüglich einer Menge zulässiger Transformationen invariant sind.

Formal

\[
F
=
[(D,R)]_{\sim}
\]

wobei

\[
\sim
\]

eine geeignete Äquivalenzrelation bezeichnet.

\end{hypothese}

Diese Sichtweise besitzt mehrere Vorteile.

\begin{itemize}

\item Form wird vollständig trägerunabhängig.

\item Unterschiedliche physikalische Realisierungen können dieselbe Form
besitzen.

\item Identität ergibt sich aus Invarianz und nicht aus Materialität.

\item Der Anschluss an moderne Mathematik (Topologie,
Kategorientheorie, Homotopie) wird wesentlich einfacher.

\end{itemize}

Diese Hypothese könnte langfristig die vorläufige Definition
\(F=(D,R,S)\) ersetzen.

%==============================================================
\section{Offene Forschungsfragen}
%==============================================================

Die DRS wirft eine Reihe grundlegender Fragen auf.

\begin{enumerate}

\item Ist Form ontologisch fundamentaler als Information?

\item Lässt sich Bestimmtheit vollständig durch Form erklären?

\item Ist Relation logisch gleichursprünglich mit Differenz?

\item Welche Minimalstruktur besitzt die einfachste mögliche Form?

\item Können Projektionsoperatoren ausschließlich zwischen Formen
abbilden?

\item Ist Bewusstsein vollständig als Klasse selbstreferenzieller Formen
beschreibbar?

\item Welche mathematische Sprache eignet sich langfristig am besten
zur Beschreibung von Formen?

\item Ist Form selbst eine emergente Eigenschaft noch fundamentalerer
Strukturen oder bildet sie den Endpunkt weiterer Reduktionen?

\end{enumerate}

%==============================================================
\section{Vorläufige Gesamtsynthese}
%==============================================================

Die in diesem Dokument entwickelte Arbeitshypothese lässt sich
gegenwärtig in folgender Struktur zusammenfassen.

\[
\boxed{
\begin{aligned}
&
\text{Differenz}
\\
\Longleftrightarrow&
\text{Relation}
\\
\Longrightarrow&
\text{Form}
\\
\Longrightarrow&
\text{Bestimmtheit}
\\
\Longrightarrow&
\text{Information}
\\
\Longrightarrow&
\text{Möglichkeiten}
\\
\Longrightarrow&
\text{Projektionsräume}
\\
\Longrightarrow&
\text{Physikalische Realität}
\end{aligned}
}
\]

Bewusstsein erscheint innerhalb dieser Sichtweise nicht als Ausnahme,
sondern als eine hochkomplexe Klasse selbstreferenzieller Formen.

%==============================================================
\section{Schnittstellen}
%==============================================================

Die DRS besitzt unmittelbare Verbindungen zu mehreren OP-Modulen.

\begin{itemize}

\item \textbf{OP:}
Die Unterscheidbarkeits-Prämisse erhält mit dem Formbegriff
eine mögliche strukturelle Präzisierung.

\item \textbf{FSO:}
Diskursdomänen können als Räume formulierbarer Formen interpretiert werden.

\item \textbf{MGO:}
Die mathematische Formalisierung von Form bildet eine natürliche
Erweiterung der dort entwickelten algebraischen und
kategorientheoretischen Grundlagen.

\item \textbf{PST:}
Projektionsoperatoren transformieren Formen zwischen verschiedenen
Beschreibungsebenen.

\item \textbf{PGT:}
Bewusstsein könnte als Klasse selbstreferenzieller Formen beschrieben
werden.

\item \textbf{DFD:}
Fragedifferentiale könnten als strukturelle Distanzen zwischen Formen
interpretiert werden.

\end{itemize}

%==============================================================
\section{Schlussbemerkung}
%==============================================================

Dieses Dokument erhebt keinen Anspruch auf eine abgeschlossene Theorie.

Es formuliert die Arbeitshypothese, dass der Begriff der Form den bislang
verwendeten Begriffen Struktur, Bestimmtheit, Relation und Information
ontologisch vorausgehen könnte.

Sollte sich diese Vermutung bestätigen, könnte die DRS zukünftig eine
vereinheitlichende Grundlage des Ontophänomenalismus bilden und als
Bindeglied zwischen OP, MGO, FSO, PST und PGT dienen.

\end{document}

